\section{Discussion}% \& Future Outlook}
\label{sec:future-outlook}

\subsection{Privacy-preserving advertising paradigms}
%
The Privacy Sandbox from Google was the most recent and extensive industrial attempt to deprecate cross-site identifiers, by relying more heavily on on-device computation and the browser to mediate access.
%
However, achieving private, but useful, ad targeting and measurement is challenging~\cite{hoganMakingSensePrivate2026}, as illustrated by the following privacy limitations uncovered in recent proposals: (a) anonymity and re-identification concerns
with FLoC~\cite{rescorlaTechnicalCommentsFLoC2021,berkePrivacyLimitationsInterestbased2022}, Topics~\cite{thomsonPrivacyAnalysisGoogles2023,jhaRobustnessTopicsAPI2023,beuginInterestdisclosingMechanismsAdvertising2024,beuginPublicReproducibleAssessment2024}, Protected Audience~\cite{thomsonProtectedAudiencePrivacy2024,longEvaluatingGooglesProtected2024}, or Apple's Private Click Measurement~\cite{thomsonAnalysisApplesPrivate2022} and Private Relay~\cite{zohaibInvestigatingTrafficAnalysis2023},
% ~\cite{rescorlaTechnicalCommentsFLoC2021,berkePrivacyLimitationsInterestbased2022,thomsonPrivacyAnalysisGoogles2023,jhaRobustnessTopicsAPI2023,beuginInterestdisclosingMechanismsAdvertising2024,beuginPublicReproducibleAssessment2024,thomsonAnalysisApplesPrivate2022,thomsonProtectedAudiencePrivacy2024,longEvaluatingGooglesProtected2024,zohaibInvestigatingTrafficAnalysis2023}
(b) flawed implementations or side-channel leakage~\cite{turatiLocalitySensitiveHashingDoes2023,aliNavigatingMurkyWaters2023,calderonioFledgingWillContinue2024,nisenoffExploitingSharedStorage2025}, and (c)  abuse~\cite{senolUnveilingImpactUserAgent2023,wieflingPrivacyMeasureTurned2024}.
%
% However, achieving private ad targeting and measurement is challenging, as illustrated by the significant privacy limitations of Google's proposals that have been uncovered.
%
% FLoC, launched in an origin trial in 2021, was withdrawn within a few months over re-identification concerns~\cite{rescorlaTechnicalCommentsFLoC2021,berkePrivacyLimitationsInterestbased2022}.
%
% It was replaced by the Topics API, whose privacy leakage was shown to worsen over time and to lead to re-identification of users as well~\cite{thomsonPrivacyAnalysisGoogles2023,jhaRobustnessTopicsAPI2023,beuginInterestdisclosingMechanismsAdvertising2024,beuginPublicReproducibleAssessment2024}.
%
% FLEDGE (later renamed Protected Audience) moved ad auctions on-device to avoid disclosing interest-group membership entirely, but auction timing and outcome side-channels were found to lead to indirect leakage~\cite{calderonioFledgingWillContinue2024,aliNavigatingMurkyWaters2023,thomsonProtectedAudiencePrivacy2024,longEvaluatingGooglesProtected2024}.
%
% Flaws in the design and browser implementation that can lead to abuse were reported as well in other Privacy Sandbox APIs~\cite{turatiLocalitySensitiveHashingDoes2023,senolUnveilingImpactUserAgent2023,wieflingPrivacyMeasureTurned2024,nisenoffExploitingSharedStorage2025,beuginLessonsAdoptionDeprecation2026}
%
% Note that Apple's Private Click Measurement (a privacy-preserving alternative for ad-click attribution alone) and Private Relay (a proxy service for iCloud users) were also criticized for providing insufficient anonymity guarantees~\cite{thomsonAnalysisApplesPrivate2022,zohaibInvestigatingTrafficAnalysis2023}

\begin{opbox}
% Every interest-disclosing replacement for third-party cookies so far has faced the same tension: the diversity that makes interests useful for advertising is what makes them re-identifying, especially under advertiser collusion. Can interest signals be designed with provable privacy-utility tradeoffs that hold even against colluding adversaries?
%% Note: there is kind of an impossibility result under the hood here
What new privacy-preserving technologies can be built by learning from issues in prior proposals? 
% How can we reliably and at scale automate the evaluation of harms, risks, and opportunities of such novel mechanisms?
\end{opbox}

By late 2025, Google deprecated most Privacy Sandbox APIs, leaving third-party cookies largely opt-in in Chrome~\cite{chavezUpdatePlansPrivacy2025,gribRiseFallGoogle2026}, while such restrictions already exist by default in Brave, Firefox, or Safari.
%
Overall, Privacy Sandbox demonstrated that besides technical challenges, economic and other (misaligned) incentives often dominate with advertising~\cite{beuginTechnicalReportNeed2025,beuginLessonsAdoptionDeprecation2026}.
%
Ideally, browsers would coordinate work on privacy-preserving advertising through the W3C's Private Advertising Technology Community Group (PATCG), chartered in 2021, rather than shipping unilateral, mutually incompatible APIs.
%
Meta and Mozilla's Interoperable Private Attribution proposal is currently under such W3C incubation.

\begin{opbox}
% Privacy Sandbox demonstrated that a tracking-relevant signal does not need to leave the device to remain useful to an advertiser, yet every browser vendor built and then abandoned its own incompatible version of this idea. 
What would a standardized, cross-browser on-device computation primitive look like that advertisers could build against regardless of which browser a user runs?
% , without requiring any single vendor to unilaterally set the privacy bar for the rest of the web?
\end{opbox}


\subsection{Tracking in other ecosystems}

While our focus was on web tracking, similar tracking mechanisms also exist in mobile apps and IoT ecosystems---using often a richer set of sensor information available via OS APIs rather than web APIs. Some key differences and similarities exist: reliance on cookies on the web and on device-level identifiers such as the Ad ID (known as \texttt{IDFA} on iOS, \texttt{AAID} on Android, or \texttt{TIFA} on Samsung Smart TVs), and vendors making different design choices across ecosystems. For instance, while Safari blocks third-party cookies by default, iOS instead asks permission to give access to device-level identifier.
% Google's Chrome and Android do not block by default third-party cookies or device-level identifiers, respectively.

\begin{opbox}
If cross-device tracking is understood theoretically, characterizing its occurrence in practice and defending against it requires more efforts. How do techniques and protections combine and diverge across ecosystems?
\end{opbox}


\subsection{AI as an emerging tracking surface}

AI is beginning to reshape web tracking on multiple fronts.
%
Browsing behaviors are shifting as AI assistants navigate the web on behalf of users. This introduces a new context where account-linked identity meets rich conversational and behavioral signals otherwise not shared on the web~\cite{vekaria2025big}.
%
As AI platforms already embed trackers~\cite{jazlan2026tracking}, introduce ads in chatbots (e.g., ChatGPT), or use LLM to infer sensitive user attributes from ordinary browsing data~\cite{chenWhenAdsBecome2026}, harms and privacy concerns may amplify.
%
Browsers have a key role to play in ensuring security and privacy of novel AI integrations.

\begin{opbox}
Will advertising and tracking shift to AI platforms, creating new tracking surfaces or amplifying the existing ones? Conversely, can real-time adaptive defenses be built with LLMs to surpass today's static heuristics?
\end{opbox}